% GENERATED FILE. Edit canonical lessons, then run npm run generate:books.
% canonical-source-hash: fnv1a64:0bc228c8c5a645c0
% canonical-lessons: PA-C355-paheli, PA-C355-bujjhna, PA-C355-mazak, PA-C355-hasauna, PA-C355-chherna

\chapter{Riddles and Jokes}
\label{ch:pa-a2-riddles-and-jokes}
\hlchaptermodality{355}

\begin{chapteropening}
\textbf{By the end of this chapter:} \emph{I can say that I like riddles and jokes.}

Complete the last lesson of chapter 355: I can say that I like riddles and jokes.
\end{chapteropening}

\section[\texorpdfstring{\pa{ਪਹੇਲੀ} (pahelī)}{ਪਹੇਲੀ (pahelī)}]{\pa{ਪਹੇਲੀ} (pahelī) — a riddle}

\label{lesson:PA-C355-paheli}



\begin{quote}
Before the new one: say the Punjabi for a queen, then the Punjabi for magic.
\end{quote}

\subsection*{You'll want to know: \pa{ਪਹੇਲੀ}}
\textbf{\pa{ਪਹੇਲੀ}} — \emph{pahelī} — \textquotedblleft{}a riddle\textquotedblright{}.

\begin{quote}
\pa{ਮੇਰੇ} \pa{ਦਾਦੇ} \pa{ਕੋਲ} \pa{ਬਹੁਤ} \pa{ਪਹੇਲੀਆਂ} \pa{ਹਨ।}
\end{quote}

\begin{quote}
\emph{mere dāde kol bahut pahelīāṁ han.}
\end{quote}

\begin{quote}
My grandfather knows lots of riddles.
\end{quote}

A feminine noun.

\begin{cousinweb}
\textbf{\pa{ਪਹੇਲੀ}} (\emph{pahelī}) is inherited from Sanskrit \emph{prahelikā}, \textquotedblleft{}riddle\textquotedblright{}; CDIAL lists Punjabi \emph{parhelī}, \textquotedblleft{}riddle\textquotedblright{}, and \emph{pahelī} under it (CDIAL 8911).
\end{cousinweb}

\begin{grammarlens}[title={What you've built}]
One more word for this chapter.
\end{grammarlens}

\subsection*{Your turn}
\begin{itemize}
\raggedright
  \item \emph{Say it:} \emph{pahelī}
  \item \emph{Say it:} \emph{pahelī}, once more
  \item \emph{Say it:} \emph{mere dāde kol bahut pahelīāṁ han.} — My grandfather knows lots of riddles.
  \item \emph{Recall:} say the Punjabi for a queen, then the Punjabi for magic, then say \emph{pahelī} again
\end{itemize}

\begin{tcolorbox}[breakable,colback=teal!4,colframe=teal!35!black,title={Before you move on}]
What does \pa{ਪਹੇਲੀ} mean? (\textquotedblleft{}A riddle\textquotedblright{}.) Say it once more.
\end{tcolorbox}
\section[\texorpdfstring{\pa{ਬੁੱਝਣਾ} (bujjhṇā)}{ਬੁੱਝਣਾ (bujjhṇā)}]{\pa{ਬੁੱਝਣਾ} (bujjhṇā) — to solve a riddle; to guess}

\label{lesson:PA-C355-bujjhna}



\begin{quote}
Before the new one: say the Punjabi for magic, then the Punjabi for a riddle.
\end{quote}

\subsection*{You'll want to know: \pa{ਬੁੱਝਣਾ}}
\textbf{\pa{ਬੁੱਝਣਾ}} — \emph{bujjhṇā} — \textquotedblleft{}to solve a riddle; to guess\textquotedblright{}.

\begin{quote}
\pa{ਮੇਰੀ} \pa{ਭੈਣ} \pa{ਹਰ} \pa{ਪਹੇਲੀ} \pa{ਬੁੱਝਦੀ} \pa{ਹੈ।}
\end{quote}

\begin{quote}
\emph{merī bhaiṇ har pahelī bujjhdī hai.}
\end{quote}

\begin{quote}
My sister solves every riddle.
\end{quote}

Past (I …): \pa{ਮੈਂ} \pa{ਪਹੇਲੀ} \pa{ਬੁੱਝੀ}, \emph{maiṁ pahelī bujjhī}.

\begin{cousinweb}
\textbf{\pa{ਬੁੱਝਣਾ}} (\emph{bujjhṇā}) is inherited from Sanskrit \emph{budhyate}, \textquotedblleft{}is awakened, understands\textquotedblright{}; CDIAL lists Punjabi \emph{bujjhṇā}, \textquotedblleft{}to guess, solve\textquotedblright{}, under it (CDIAL 9279).
\end{cousinweb}

\begin{grammarlens}[title={What you've built}]
One more word for this chapter.
\end{grammarlens}

\subsection*{Your turn}
\begin{itemize}
\raggedright
  \item \emph{Say it:} \emph{bujjhṇā}
  \item \emph{Say it:} \emph{bujjhṇā}, once more
  \item \emph{Say it:} \emph{merī bhaiṇ har pahelī bujjhdī hai.} — My sister solves every riddle.
  \item \emph{Recall:} say the Punjabi for magic, then the Punjabi for a riddle, then say \emph{bujjhṇā} again
  \item \emph{Write it:} copy \textbf{\pa{ਬੁੱਝਣਾ}} once, checking the whole word
\end{itemize}

\begin{tcolorbox}[breakable,colback=teal!4,colframe=teal!35!black,title={Before you move on}]
What does \pa{ਬੁੱਝਣਾ} mean? (\textquotedblleft{}To solve a riddle; to guess\textquotedblright{}.) Say it once more.
\end{tcolorbox}
\section[\texorpdfstring{\pa{ਮਜ਼ਾਕ} (mazāk)}{ਮਜ਼ਾਕ (mazāk)}]{\pa{ਮਜ਼ਾਕ} (mazāk) — a joke; joking}

\label{lesson:PA-C355-mazak}



\begin{quote}
Before the new one: say the Punjabi for a riddle, then the Punjabi for to solve a riddle.
\end{quote}

\subsection*{You'll want to know: \pa{ਮਜ਼ਾਕ}}
\textbf{\pa{ਮਜ਼ਾਕ}} — \emph{mazāk} — \textquotedblleft{}a joke; joking\textquotedblright{}.

\begin{quote}
\pa{ਮੇਰਾ} \pa{ਦੋਸਤ} \pa{ਬਹੁਤ} \pa{ਮਜ਼ਾਕ} \pa{ਕਰਦਾ} \pa{ਹੈ।}
\end{quote}

\begin{quote}
\emph{merā dost bahut mazāk kardā hai.}
\end{quote}

\begin{quote}
My friend jokes a lot.
\end{quote}

A masculine noun.

\begin{cousinweb}
\textbf{\pa{ਮਜ਼ਾਕ}} (\emph{mazāk}) is a loan, usually traced through Persian to Arabic \emph{madhāq}, \textquotedblleft{}taste\textquotedblright{}.
\end{cousinweb}

\begin{grammarlens}[title={What you've built}]
One more word for this chapter.
\end{grammarlens}

\subsection*{Your turn}
\begin{itemize}
\raggedright
  \item \emph{Say it:} \emph{mazāk}
  \item \emph{Say it:} \emph{mazāk}, once more
  \item \emph{Say it:} \emph{merā dost bahut mazāk kardā hai.} — My friend jokes a lot.
  \item \emph{Recall:} say the Punjabi for a riddle, then the Punjabi for to solve a riddle, then say \emph{mazāk} again
  \item \emph{Write it:} copy \textbf{\pa{ਮਜ਼ਾਕ}} once, checking the whole word
\end{itemize}

\begin{tcolorbox}[breakable,colback=teal!4,colframe=teal!35!black,title={Before you move on}]
What does \pa{ਮਜ਼ਾਕ} mean? (\textquotedblleft{}A joke; joking\textquotedblright{}.) Say it once more.
\end{tcolorbox}
\section[\texorpdfstring{\pa{ਹਸਾਉਣਾ} (hasāuṇā)}{ਹਸਾਉਣਾ (hasāuṇā)}]{\pa{ਹਸਾਉਣਾ} (hasāuṇā) — to make someone laugh}

\label{lesson:PA-C355-hasauna}



\begin{quote}
Before the new one: say the Punjabi for to solve a riddle, then the Punjabi for a joke.
\end{quote}

\subsection*{You'll want to know: \pa{ਹਸਾਉਣਾ}}
\textbf{\pa{ਹਸਾਉਣਾ}} — \emph{hasāuṇā} — \textquotedblleft{}to make someone laugh\textquotedblright{}.

\begin{quote}
\pa{ਅਮਨ} \pa{ਸਭ} \pa{ਨੂੰ} \pa{ਹਸਾਉਂਦਾ} \pa{ਹੈ।}
\end{quote}

\begin{quote}
\emph{Aman sabh nū̃ hasāundā hai.}
\end{quote}

\begin{quote}
Aman makes everyone laugh.
\end{quote}

Past (I …): \pa{ਮੈਂ} \pa{ਬੱਚਿਆਂ} \pa{ਨੂੰ} \pa{ਹਸਾਇਆ}, \emph{maiṁ bachchiāṁ nū̃ hasāiā}.

\begin{cousinweb}
\textbf{\pa{ਹਸਾਉਣਾ}} (\emph{hasāuṇā}) is the causative partner of \emph{hassṇā}, \textquotedblleft{}to laugh\textquotedblright{}, which goes back to a reconstructed \emph{*hasyate}, \textquotedblleft{}one laughs\textquotedblright{}; CDIAL lists Punjabi \emph{hassṇā} under it (CDIAL 14048).
\end{cousinweb}

\begin{grammarlens}[title={What you've built}]
One more word for this chapter.
\end{grammarlens}

\subsection*{Your turn}
\begin{itemize}
\raggedright
  \item \emph{Say it:} \emph{hasāuṇā}
  \item \emph{Say it:} \emph{hasāuṇā}, once more
  \item \emph{Say it:} \emph{Aman sabh nū̃ hasāundā hai.} — Aman makes everyone laugh.
  \item \emph{Recall:} say the Punjabi for to solve a riddle, then the Punjabi for a joke, then say \emph{hasāuṇā} again
  \item \emph{Write it:} copy \textbf{\pa{ਹਸਾਉਣਾ}} once, checking the whole word
\end{itemize}

\begin{tcolorbox}[breakable,colback=teal!4,colframe=teal!35!black,title={Before you move on}]
What does \pa{ਹਸਾਉਣਾ} mean? (\textquotedblleft{}To make someone laugh\textquotedblright{}.) Say it once more.
\end{tcolorbox}
\section[\texorpdfstring{\pa{ਛੇੜਨਾ} (chheṛnā)}{ਛੇੜਨਾ (chheṛnā)}]{\pa{ਛੇੜਨਾ} (chheṛnā) — to tease}

\label{lesson:PA-C355-chherna}



\begin{quote}
Before the new one: say the Punjabi for a joke, then the Punjabi for to make someone laugh.
\end{quote}

\subsection*{You'll want to know: \pa{ਛੇੜਨਾ}}
\textbf{\pa{ਛੇੜਨਾ}} — \emph{chheṛnā} — \textquotedblleft{}to tease\textquotedblright{}.

\begin{quote}
\pa{ਮੇਰਾ} \pa{ਭਰਾ} \pa{ਮੈਨੂੰ} \pa{ਹਰ} \pa{ਰੋਜ਼} \pa{ਛੇੜਦਾ} \pa{ਹੈ।}
\end{quote}

\begin{quote}
\emph{merā bharā mainū̃ har roz chheṛdā hai.}
\end{quote}

\begin{quote}
My brother teases me every day.
\end{quote}

Past (I …): \pa{ਮੈਂ} \pa{ਭੈਣ} \pa{ਨੂੰ} \pa{ਛੇੜਿਆ}, \emph{maiṁ bhaiṇ nū̃ chheṛiā}.

\begin{cousinweb}
\textbf{\pa{ਛੇੜਨਾ}} (\emph{chheṛnā}) is inherited from a reconstructed \emph{*ciḍ}, \textquotedblleft{}be angry\textquotedblright{}, with \emph{*cheḍ}, \textquotedblleft{}make angry\textquotedblright{}; CDIAL lists Punjabi \emph{cheṛnā}, \textquotedblleft{}to provoke\textquotedblright{}, under it (CDIAL 4794).
\end{cousinweb}

\begin{grammarlens}[title={What you've built}]
One more word for this chapter.
\end{grammarlens}

\subsection*{Your turn}
\begin{itemize}
\raggedright
  \item \emph{Say it:} \emph{chheṛnā}
  \item \emph{Say it:} \emph{chheṛnā}, once more
  \item \emph{Say it:} \emph{merā bharā mainū̃ har roz chheṛdā hai.} — My brother teases me every day.
  \item \emph{Recall:} say the Punjabi for a joke, then the Punjabi for to make someone laugh, then say \emph{chheṛnā} again
  \item \emph{Write it:} copy \textbf{\pa{ਛੇੜਨਾ}} once, checking the whole word
\end{itemize}

\begin{tcolorbox}[breakable,colback=teal!4,colframe=teal!35!black,title={Before you move on}]
What does \pa{ਛੇੜਨਾ} mean? (\textquotedblleft{}To tease\textquotedblright{}.) Say it once more.
\end{tcolorbox}
